%% Template Jinja2+LaTeX. Variables: estudiante, contexto/muestra/sol{1,2,3}_*, modo, primero, ultimo.
%% El bloque "primero" emite el preámbulo una sola vez; "ultimo" cierra \end{document}.
{% if primero -%}
\documentclass[10pt]{article}
\usepackage[spanish]{babel}
\usepackage[utf8]{inputenc}
\usepackage{amsmath}    %% \dfrac, \bar, ecuaciones
\usepackage{longtable, array, pdflscape}  %% hoja de respuestas al inicio de las soluciones
\usepackage{fancyhdr}   %% header/footer personalizados
\usepackage{siunitx}
\usepackage{lastpage}   %% \pageref{LastPage} para "X de N"
\usepackage{geometry}
\usepackage{xcolor}
\usepackage{tikz}
\usepackage{pgfplots}   %% gráfico de barras pregunta 3
\usepackage{enumitem}   %% enumerate compacto en soluciones

\pgfplotsset{compat=1.18}
\geometry{top=1.5cm, bottom=1.5cm, left=1.2cm, right=1.2cm}

\pagestyle{fancy}
\fancyhf{}
\lhead{\small ESTAD\'ISTICA}                       %% \'I = Í (acento sobre la I)
\rhead{\small \today}
\renewcommand{\headrulewidth}{0.5pt}
\rfoot{\thepage\ de \pageref{LastPage}}            %% requiere 2 pasadas de pdflatex
{% if modo == 'soluciones' %}\lfoot{\scriptsize Semilla: {{ semilla }}}{% endif %}  %% permite repetir la corrida
\setlength{\headheight}{15pt}
\setlength{\parindent}{0pt}

\begin{document}
{{ clave_tex }}          %% vacío en las prácticas; en las soluciones, la hoja de respuestas
{%- endif %}

\label{inicio_{{ estudiante_id }}}                 %% etiqueta única por estudiante
\begin{center}
    {\Large \textbf{Práctica de Estadística Descriptiva}}\\[0.3em]
    {\normalsize Estudiante: \textbf{\underline{ {{ estudiante }} }} }
\end{center}

\vspace{0.15cm}

\noindent\textbf{Indicador:} Caracterizar un grupo de datos mediante medidas estadísticas de resumen: moda, media aritmética, mediana, máximo, mínimo y recorrido; además de interpretar su representación gráfica.

\vspace{0.1cm}

%% Caja gris con la indicación general (mismo patrón se reutiliza para las muestras).
\noindent\fcolorbox{black}{gray!10}{\parbox{\dimexpr\textwidth-2\fboxsep-2\fboxrule\relax}{
\small\centering
\textit{Recuerde mostrar todos sus procedimientos. Puede usar calculadora para operaciones aritméticas básicas.}
}}

\vspace{0.1cm}
\hrule
\vspace{0.15cm}

{% if modo == 'practicas' %}                       %% ---------- ENUNCIADOS ----------

%% PREGUNTA 1: estadísticos básicos sobre muestra ordenada uniforme.
\textbf{1. } Se registraron los siguientes datos sobre {{ contexto1 }}:

\noindent\fcolorbox{black}{gray!10}{\parbox{\dimexpr\linewidth-2\fboxsep-2\fboxrule\relax}{
\small\centering
{{ muestra1 }}
}}

Determine las medidas de resumen. Primero ordene los datos y luego calcule:
\begin{enumerate}
\item[\small a)] La moda (dato que más se repite, si existe)
\item[\small b)] El mínimo y el máximo
\item[\small c)] La media aritmética aplicando la fórmula:
$$\bar{x} = \frac{\text{Suma de todos los datos}}{\text{Cantidad de datos}}$$
\item[\small d)] El recorrido: $R = M\acute{a}x - M\acute{\imath}n$
\item[\small e)] La mediana (dato central de los datos ordenados; con una cantidad par, promedio de los dos centrales)
\end{enumerate}

\vspace{0.25cm}

%% PREGUNTA 2: la muestra trae valores repetidos para forzar moda real.
\textbf{2. } Analice el siguiente conjunto de datos de {{ contexto2 }}:

\noindent\fcolorbox{black}{gray!10}{\parbox{\dimexpr\linewidth-2\fboxsep-2\fboxrule\relax}{
\small\centering
{{ muestra2 }}
}}

\begin{enumerate}
\item[\small a)] Organice los datos de menor a mayor
\item[\small b)] Determine la moda (el valor con mayor frecuencia)
\item[\small c)] Calcule la media aritmética mostrando la suma y la división
\item[\small d)] Calcule el recorrido e interprete qué significa este valor
\item[\small e)] Determine la mediana a partir de los datos ordenados
\end{enumerate}

\vspace{0.25cm}

%% PREGUNTA 3: gráfico de barras + línea de la media (TikZ/pgfplots).
\textbf{3. } Análisis gráfico de los datos de {{ contexto3 }}:

\noindent\fcolorbox{black}{gray!10}{\parbox{\dimexpr\linewidth-2\fboxsep-2\fboxrule\relax}{
\small\centering
{{ muestra3 }}
}}

\vspace{0.15cm}

\noindent\textbf{Gráfico de Barras:}

\vspace{0.1cm}

\begin{tikzpicture}[scale=0.9]
\begin{axis}[
    title style={font=\Large\bfseries},
    title={Representación de Datos},
    xlabel={Posición}, ylabel={Valor},
    ymin=0, ymax={{ (sol3_maximo * 1.15)|int }},  %% 15% de aire arriba
    xtick={1,...,{{ sol3_datos_ordenados|length }}},
    axis lines=left, width=10cm, height=6cm,
    legend style={draw=none, fill=none}, legend entries={},
]
\addplot[ybar, bar shift=0pt, fill=blue!70, draw=black, line width=0.75pt] coordinates {
{% for val in sol3_datos_ordenados -%}({{ loop.index }}, {{ val }}){% endfor -%}};
\addplot[red, dashed, thick, mark=none] coordinates {  %% línea horizontal en la media
(0.5, {{ "%.2f"|format(sol3_media) }}) ({{ sol3_datos_ordenados|length + 0.5 }}, {{ "%.2f"|format(sol3_media) }})};
\end{axis}
\end{tikzpicture}

\vspace{0.15cm}

\begin{enumerate}
\item[\small a)] Identifique el mínimo y el máximo en el gráfico
\item[\small b)] Localice dónde está la media
\item[\small c)] Calcule la mediana y compárela con la media. ¿Qué indica la diferencia entre ambas?
\item[\small d)] Escriba una breve interpretación: ¿qué le dice este gráfico sobre los datos de {{ contexto3 }}?
\end{enumerate}

\vspace{0.25cm}

{% else %}                                         %% ---------- SOLUCIONES (compactas) ----------

%% Reduce el espaciado de fórmulas para que las 3 soluciones quepan en una página.
\begingroup
\setlength{\abovedisplayskip}{6pt}
\setlength{\belowdisplayskip}{6pt}
\setlength{\abovedisplayshortskip}{6pt}
\setlength{\belowdisplayshortskip}{6pt}

\vspace{0.3cm}

\textbf{1. {{ contexto1 }}} \par
\textit{Datos:} {{ muestra1 }}
\begin{enumerate}[leftmargin=1.4em, itemsep=3pt, topsep=4pt, parsep=0pt]
\item[a)] \textbf{Moda:} {{ sol1_moda_texto }}
\item[b)] \textbf{Mínimo:} {{ sol1_minimo }} \quad \textbf{Máximo:} {{ sol1_maximo }}
\item[c)] \textbf{Media:} $\bar{x} = \dfrac{ {{ sol1_datos_ordenados|join('+') }} }{ {{sol1_n}} } = \dfrac{ {{sol1_suma}} }{ {{sol1_n}} } = {{ "%.2f"|format(sol1_media) }}$
\item[d)] \textbf{Recorrido:} $R = {{ sol1_maximo }} - {{ sol1_minimo }} = {{ sol1_recorrido }}$
\item[e)] \textbf{Mediana:} $Me = {{ sol1_mediana_tex }}$
\end{enumerate}

\vspace{0.3cm}

\textbf{2. {{ contexto2 }}} \par
\textit{Datos:} {{ muestra2 }}
\begin{enumerate}[leftmargin=1.4em, itemsep=3pt, topsep=4pt, parsep=0pt]
\item[a)] \textbf{Datos ordenados:} {{ sol2_datos_ordenados|join(', ') }}
\item[b)] \textbf{Moda:} {{ sol2_moda_texto }}
\item[c)] \textbf{Media:} $\bar{x} = \dfrac{ {{ sol2_datos_ordenados|join('+') }} }{ {{sol2_n}} } = \dfrac{ {{sol2_suma}} }{ {{sol2_n}} } = {{ "%.2f"|format(sol2_media) }}$
\item[d)] \textbf{Recorrido:} $R = {{ sol2_maximo }} - {{ sol2_minimo }} = {{ sol2_recorrido }}$ (rango entre el dato menor y el mayor)
\item[e)] \textbf{Mediana:} $Me = {{ sol2_mediana_tex }}$
\end{enumerate}

\vspace{0.3cm}

\textbf{3. {{ contexto3 }} -- Análisis gráfico} \par
\textit{Datos:} {{ muestra3 }}

%% Mismo gráfico que la pregunta 3 del enunciado, pero más pequeño para dejar espacio a los incisos.
\begin{center}
\begin{tikzpicture}[scale=0.85]
\begin{axis}[
    title style={font=\normalsize\bfseries},
    title={Representación de Datos},
    xlabel={Posición}, ylabel={Valor},
    label style={font=\small}, tick label style={font=\small},
    ymin=0, ymax={{ (sol3_maximo * 1.15)|int }},
    xtick={1,...,{{ sol3_datos_ordenados|length }}},
    axis lines=left, width=11cm, height=5cm,
    legend style={draw=none, fill=none}, legend entries={},
]
\addplot[ybar, bar shift=0pt, fill=blue!70, draw=black, line width=0.7pt] coordinates {
{% for val in sol3_datos_ordenados -%}({{ loop.index }}, {{ val }}){% endfor -%}};
\addplot[red, dashed, thick, mark=none] coordinates {
(0.5, {{ "%.2f"|format(sol3_media) }}) ({{ sol3_datos_ordenados|length + 0.5 }}, {{ "%.2f"|format(sol3_media) }})};
\end{axis}
\end{tikzpicture}
\end{center}

\begin{enumerate}[leftmargin=1.4em, itemsep=3pt, topsep=2pt, parsep=0pt]
\item[a)] \textbf{Mínimo:} {{ sol3_minimo }} \quad \textbf{Máximo:} {{ sol3_maximo }} (extremos en el gráfico)
\item[b)] \textbf{Media:} {{ "%.2f"|format(sol3_media)|replace('.', ',') }} (línea roja punteada)
\item[c)] \textbf{Mediana:} $Me = {{ sol3_mediana_tex }}$
\item[d)] \textit{Interpretación:} {{ sol3_interpretacion }}
\end{enumerate}

\endgroup

{% endif %}

\label{fin_{{ estudiante_id }}}
{% if not ultimo %}\newpage{% endif %}             %% salto entre estudiantes, no después del último
{% if ultimo %}
\end{document}
{% endif %}
